% Options for packages loaded elsewhere
\PassOptionsToPackage{unicode}{hyperref}
\PassOptionsToPackage{hyphens}{url}
\PassOptionsToPackage{space}{xeCJK}
%
\documentclass[
]{ctexart}
\usepackage{amsmath,amssymb}
\usepackage{iftex}
\ifPDFTeX
  \usepackage[T1]{fontenc}
  \usepackage[utf8]{inputenc}
  \usepackage{textcomp} % provide euro and other symbols
\else % if luatex or xetex
  \usepackage{unicode-math} % this also loads fontspec
  \defaultfontfeatures{Scale=MatchLowercase}
  \defaultfontfeatures[\rmfamily]{Ligatures=TeX,Scale=1}
\fi
\usepackage{lmodern}
\ifPDFTeX\else
  % xetex/luatex font selection
  \ifXeTeX
    \usepackage{xeCJK}
    \setCJKmainfont[]{Songti SC}
          \fi
  \ifLuaTeX
    \usepackage[]{luatexja-fontspec}
    \setmainjfont[]{Songti SC}
  \fi
\fi
% Use upquote if available, for straight quotes in verbatim environments
\IfFileExists{upquote.sty}{\usepackage{upquote}}{}
\IfFileExists{microtype.sty}{% use microtype if available
  \usepackage[]{microtype}
  \UseMicrotypeSet[protrusion]{basicmath} % disable protrusion for tt fonts
}{}
\makeatletter
\@ifundefined{KOMAClassName}{% if non-KOMA class
  \IfFileExists{parskip.sty}{%
    \usepackage{parskip}
  }{% else
    \setlength{\parindent}{0pt}
    \setlength{\parskip}{6pt plus 2pt minus 1pt}}
}{% if KOMA class
  \KOMAoptions{parskip=half}}
\makeatother
\usepackage{xcolor}
\usepackage[margin=2.2cm]{geometry}
\setlength{\emergencystretch}{3em} % prevent overfull lines
\providecommand{\tightlist}{%
  \setlength{\itemsep}{0pt}\setlength{\parskip}{0pt}}
\setcounter{secnumdepth}{5}
\renewcommand{\and}{\\}
\ifLuaTeX
  \usepackage{selnolig}  % disable illegal ligatures
\fi
\IfFileExists{bookmark.sty}{\usepackage{bookmark}}{\usepackage{hyperref}}
\IfFileExists{xurl.sty}{\usepackage{xurl}}{} % add URL line breaks if available
\urlstyle{same}
\hypersetup{
  pdftitle={Clustering},
  pdfauthor={Rui},
  hidelinks,
  pdfcreator={LaTeX via pandoc}}

\title{Clustering}
\author{Rui}
\date{\texttt{r\ format(Sys.Date(),\ \textquotesingle{}\%B\ \%d,\ \%Y\textquotesingle{})}}

\begin{document}
\maketitle

\renewcommand*\contentsname{目录}
{
\setcounter{tocdepth}{3}
\tableofcontents
}
\texttt{\{r\ results=\textquotesingle{}hide\textquotesingle{},message=FALSE\}\ library("parameters")\ \ \#for\ clustering\ summary\ library("see")\ library("factoextra")\ \ \#for\ clustering\ plotting\ library(bruceR)\ library(tidyverse)}

\texttt{\{r\}\ athlete\ =\ import("data13-02.xlsx")\ \%\textgreater{}\%\ select(-1)}

\hypertarget{k-means}{%
\section{K-Means}\label{k-means}}

\hypertarget{computation}{%
\subsection{Computation}\label{computation}}

\begin{itemize}
\tightlist
\item
  \texttt{x}: numeric matrix-like data
\item
  \texttt{centers}: either the number of clusters, say kk, or a set of
  initial (distinct) cluster centres. If a number, a random set of
  (distinct) rows in x is chosen as the initial centres.
\item
  \texttt{iter.max}: the maximum number of iterations allowed.
\item
  \texttt{nstart}: if centers is a number, how many random sets should
  be chosen. Trying several random starts is often recommended.
\end{itemize}

\texttt{\{r\}\ kmean\_results\ =\ kmeans(athlete,\ centers\ =\ 4,\ \ \ \ \ \ \ \ \ \ \ \ \ \ \ \ \ \ \ \ \ \ \ \ \ iter.max\ =\ 10,\ nstart\ =\ 1)}

kmeans returns an object of class ``kmeans'' which has a print and a
fitted method. It is a list with at least the following components:

\begin{itemize}
\tightlist
\item
  \texttt{cluster} : A vector of integers (from 1:k) indicating the
  cluster to which each point is allocated.
\item
  \texttt{centers}: A matrix of cluster centres.
\item
  \texttt{totss}: The total sum of squares.
\item
  \texttt{withinss}: Vector of within-cluster sum of squares, one
  component per cluster.
\item
  \texttt{tot.withinss} :Total within-cluster sum of squares,
  i.e.~sum(withinss).
\item
  \texttt{betweenss}: The between-cluster sum of squares,
  i.e.~\texttt{totss}-\texttt{tot.withinss}.
\item
  \texttt{size}: The number of points in each cluster.
\end{itemize}

\texttt{\{r\}\ kmean\_results\$cluster\ \ \#group\ labels\ for\ each\ observation}

\texttt{\{r\}\ kmean\_results\$centers\ \ \#centers}

\texttt{\{r\ results=\textquotesingle{}hold\textquotesingle{}\}\ kmean\_results\$totss\ \ \#\ total\ variance\ kmean\_results\$withinss\ \ \#\ within-group\ variance\ kmean\_results\$tot.withinss\ \#\ sum\ of\ within\ SS\ kmean\_results\$betweenss\ \ \#between-group\ variance}

\texttt{\{r\}\ kmean\_results\$size\ \ \#\ how\ many\ cases\ in\ each\ cluster}

\hypertarget{visualization}{%
\subsection{Visualization}\label{visualization}}

Use \texttt{factoextra::fviz\_cluster(object,\ data)}. For more details
about useful parameters, see the help document. Remember that
\texttt{object} and \texttt{data} are two most indispensable parameters.

\texttt{\{r\}\ fviz\_cluster(kmean\_results,\ data\ =\ athlete)}

One way to start is from random centers; another way is to specify a few
sets of centers and start from them.

Clearly, starting from different centers, either random or specified,
the results may vary, though not significant. If the iteration will
converge eventually, then given sufficiently-large iteration time, the
center will be determined and fixed.

\texttt{\{r\}\ initial\_centers\ =\ import("data13-02a.xlsx")\ \%\textgreater{}\%\ select(-1)\ kmean\_results\_initialized\ =\ kmeans(athlete,\ initial\_centers\ ,\ \ \ \ \ \ \ \ \ \ \ \ \ \ \ \ \ \ \ \ \ \ \ \ \ \ \ \ \ \ \ \ \ \ \ \ \ iter.max\ =\ 10,\ nstart\ =\ 1)\ \ fviz\_cluster(kmean\_results\_initialized,\ data\ =\ athlete)}

\hypertarget{hierachical-clustering}{%
\section{Hierachical Clustering}\label{hierachical-clustering}}

\texttt{\{r\}\ automobile\ =\ import("data13-01.xlsx")\ \%\textgreater{}\%\ \ \ \ select(c(\textquotesingle{}type\textquotesingle{},\textquotesingle{}price\textquotesingle{},\textquotesingle{}engine\_s\textquotesingle{},\textquotesingle{}horsepow\textquotesingle{},\ \ \ \ \ \ \ \ \ \ \ \ \textquotesingle{}wheelbas\textquotesingle{},\textquotesingle{}width\textquotesingle{},\textquotesingle{}length\textquotesingle{},\textquotesingle{}curb\_wgt\textquotesingle{},\ \ \ \ \ \ \ \ \ \ \ \ \textquotesingle{}fuel\_cap\textquotesingle{},\textquotesingle{}mpg\textquotesingle{}))\ \%\textgreater{}\%\ drop\_na()\ \%\textgreater{}\%\ scale()}

Use ``\texttt{hclust()}'' from base R and transmit a
\textbf{dissimilarity} structure produced by ``\texttt{dist()}''.

\texttt{\{r\}\ distMat\ =\ dist(automobile,\ method\ =\ "euclidean")\ \#\ Euclidean\ distance\ matrix\ hc\ =\ hclust(distMat)\ \ \#\ hierarchical\ clustering,\ using\ the\ default\ method}

Use ``\texttt{plot(hclust\_obj)}'' to visualize the hierarchical
clustering. Note that you should specify ``\texttt{hang\ =\ -1}'' to
ensure all nodes touch the same horizontal line and looks decent. If the
sample isn't that large, set ``\texttt{labels\ =\ TRUE}'' (by default)
to see the corresponding observations.

\texttt{\{r\}\ plot(hc,labels=FALSE,hang\ =\ -1)\ plot(hc,labels=FALSE)}

Moreover, you can draw dendrogram with grouping.

\texttt{\{r\ warning=FALSE\}\ fviz\_dend(hc,k=3,\ rect\ =\ T)}

Also, you can also check how many members for each cluster and make a
frequency table.

Note that the grouping might change if we change cluster number. And
this in essence is the pros of hierarchical clustering.

\texttt{\{r\}\ memberLabels\ =\ cutree(hc,3)\ table(memberLabels)}

Then, we can get the centers of cluters by ``\texttt{mean}'' summary
statistics. Use ``\texttt{aggregate()}'' to achieve that.

\begin{itemize}
\tightlist
\item
  \texttt{x}: usually a data frame.
\item
  \texttt{by}: a list of grouping elements, each \textbf{as long} as the
  variables in the data frame x. The elements are coerced to factors
  before use.
\item
  \texttt{FUN}: a function to compute the summary statistics which can
  be applied to all data subsets.
\end{itemize}

\texttt{\{r\}\ aggregate(automobile,list(memberLabels),mean)}

Moreover, we can use Silhouette plot to check whether some cases are
outliers.

\texttt{\{r\}\ library(cluster)\ plot(silhouette(cutree(hc,3),\ distMat))}

\hypertarget{no.-of-clusters}{%
\section{No.~of Clusters}\label{no.-of-clusters}}

\hypertarget{within-cluster-ss}{%
\subsection{Within-Cluster SS}\label{within-cluster-ss}}

Looking for the ``elbow'' where the SS starts to slow down its drop.

\texttt{\{r\}\ fviz\_nbclust(automobile,\ kmeans,\ method\ =\ "wss",\ print.summary=TRUE)}

here suggests 2 or 3.

\hypertarget{silhouette-method}{%
\subsection{Silhouette Method}\label{silhouette-method}}

Looking for the value that sets the highest average silhouette width,
which indicates the optimal number of clusters.

\texttt{\{r\}\ fviz\_nbclust(automobile,\ kmeans,\ method\ =\ "silhouette",print.summary=TRUE)}

\hypertarget{majority-rule}{%
\subsection{Majority Rule}\label{majority-rule}}

\texttt{n\_clusters()} is the main function to find out the optimal
numbers of clusters present in the data based on the maximum consensus
of a large number of methods.

Essentially, there exist many methods to determine the optimal number of
clusters, each with pros and cons, benefits and limitations. The main
\texttt{n\_clusters()} function proposes to run all of them, and find
out the number of clusters that is suggested by the \textbf{majority} of
methods (in case of ties, it will select the most parsimonious solution
with fewer clusters).

Note that we also implement some specific, commonly used methods, like
the Elbow or the Gap method, with their own visualization
functionalities.

Also note that the \texttt{n\_clusters()} function will do the
normalization automatically (by default).

\texttt{\{r\ results=\textquotesingle{}asis\textquotesingle{}\}\ n\_clusters(data.frame(automobile))\ plot(n\_clusters(data.frame(automobile)))}

\hypertarget{two-step-clustering}{%
\section{Two Step Clustering}\label{two-step-clustering}}

Some people prefer using hierarchical methods first to obtain the number
of clusters and their centroids. Then use these as inputs to run
non-hierarchical analysis. This approach is called Two-Step Clustering.

\texttt{\{r\ eval=FALSE\}\ \#\ two\ step\ clustering\ in\ R,\ not\ required\ library(prcr)\ d\ \textless{}-\ pisaUSA15\ m3\ \textless{}-\ create\_profiles\_cluster(d,\ \ \ \ \ \ \ \ \ \ \ \ \ \ \ \ \ \ \ \ \ \ \ \ \ \ \ \ \ \ \ broad\_interest,\ enjoyment,\ instrumental\_mot,\ self\_efficacy,\ \ \ \ \ \ \ \ \ \ \ \ \ \ \ \ \ \ \ \ \ \ \ \ \ \ \ \ \ \ \ n\_profiles\ =\ 3)\ summary(m3)}

\hypertarget{clustering-for-cases-and-variables}{%
\section{Clustering for Cases and
Variables}\label{clustering-for-cases-and-variables}}

Infrastructure construction includes various dimensions such as
transportation, postal and telecommunications, water supply and power
supply. ``construction.xlsx'' records the development level of 16
infrastructure construction aspects in 7 countries. Please answer the
following questions:

\begin{enumerate}
\def\labelenumi{\arabic{enumi}.}
\item
  Please use hierarchical methods to determine how many different types
  of infrastructure construction are included in these 16 indicators,
  and briefly explain the possible meanings of each type.
\item
  Calculate the average value of each type of indicator for each country
  in the previous question as the overall measure of infrastructure
  construction for that type (assuming this operation is reasonable).
  Use the k-means method to cluster the 7 countries based on their
  different overall indicators, distinguish between developed and
  developing countries, and determine which countries belong to each
  category.
\end{enumerate}

\hypertarget{q1}{%
\subsection{Q1}\label{q1}}

\texttt{\{r\ results=\textquotesingle{}hide\textquotesingle{},message=FALSE\}\ library(factoextra)\ library(showtext)\ showtext\_auto()}

\texttt{\{r\ message=FALSE,warning=FALSE\}\ cons\ =\ import("construction.xlsx")\ rownames(cons)\ =\ cons{[},1{]}\ cons\ =\ cons\ \%\textgreater{}\%\ select(-1)\ attach(cons)}

\texttt{\{r\}\ distance\ =\ dist(cons)\ hc\ =\ hclust(distance)}

\texttt{\{r\ warning=FALSE,message=FALSE\}\ plot(hc,hang\ =\ -1)}

From the cluster dendrogram above, it's suggested to classify all 16
indexes into 3 grand categories. From indexes contained in each
category, we can infer that the three dimensions are:

\begin{itemize}
\tightlist
\item
  Medical resources
\item
  Education level
\item
  Traffic capacity
\end{itemize}

\hypertarget{q2}{%
\subsection{Q2}\label{q2}}

\texttt{\{r\}\ grouplabel\ =\ cutree(hc,k\ =\ 3)}

\texttt{\{r\}\ km\_country\ =\ aggregate(cons,list(grouplabel),mean)\ \%\textgreater{}\%\ \ \ \ select(-1)\ \%\textgreater{}\%\ t()\ \%\textgreater{}\%\ \ \ \ kmeans(centers\ =\ 2)\ country\_label\ =\ km\_country\$cluster\ country\_label}

From the \(k\)-means method, we can see that the seven countries are
classified into two grand categories, developed and developing
countries, where the first group consists of D, E \& F, and all else as
the second group. However, clustering analysis doesn't directly tell us
the attributes of two groups, therefore, we have to ``label'' the groups
by checking the data.

Note that here D, E \& F are labeled as group ``2''.

\texttt{\{r\ results=\textquotesingle{}asis\textquotesingle{}\}\ cons\ \%\textgreater{}\%\ t()\ \%\textgreater{}\%\ as.data.frame()\ \%\textgreater{}\%\ mutate(\ \ \ develop\ =\ country\_label\ )\ \%\textgreater{}\%\ group\_by(develop)\ \%\textgreater{}\%\ summarise\_all(mean)\ \%\textgreater{}\%\ \ \ \ stargazer::stargazer(header\ =\ FALSE,\ \ \ \ \ \ \ \ \ \ \ \ \ \ \ \ \ \ \ \ \ \ \ \ title\ =\ \textquotesingle{}Index\ Comparison\ Between\ Groups\textquotesingle{},\ \ \ \ \ \ \ \ \ \ \ \ \ \ \ \ \ \ \ \ \ \ \ \ flip\ =\ T,summary\ =\ FALSE)}

From the summarization table, apparently that among all the 16 indexes,
group 1 outbeats group 2 (i.e., D, E \& F). Therefore, we can say

\begin{itemize}
\tightlist
\item
  Developing countries: D, E, F;
\item
  Developed countries: A, B, C, G.
\end{itemize}

\texttt{\{r\}\ fviz\_cluster(km\_country,\ \ \ \ \ \ \ \ \ \ \ \ \ \ data\ =\ aggregate(cons,list(grouplabel),mean)\ \%\textgreater{}\%\ \ \ \ \ \ \ \ \ \ \ \ \ \ \ \ \ select(-1)\ \%\textgreater{}\%\ t())}

All done. Finally, \texttt{\{r\}\ detach(cons)}

\end{document}
